% Solutions of 18-heaps.
\section*{Chapter \ref{ch:b1:heaps} --- Heaps and Priority Queues}

\begin{solution}{exo:b1:heaps:1}
$[5]$, $[3, 5]$, $[3, 5, 8]$; 1 enters at index 3, climbs past 5 and 3: $[1, 3, 8, 5]$; 4 enters at index 4 below 3 and
stays: $[1, 3, 8, 5, 4]$; 2 enters at index 5 below 8, climbs once and stops below 1: $[1, 3, 2, 5, 4, 8]$.
\end{solution}

\begin{solution}{exo:b1:heaps:2}
The root 1 is returned and the last \omterm{def:b1:induction-and-termination:rooted-tree}{leaf} 8 moves to the root: $[8, 3, 2, 5, 4]$. Its smaller child is 2 (index 2): $[2,
3, 8, 5, 4]$; index 2 has no child in a heap of 5. Result $[2, 3, 8, 5, 4]$.
\end{solution}

\begin{solution}{exo:b1:heaps:3}
Parent $\lfloor 12/2 \rfloor = 6$, children 27 and 28. Depth 3 holds the indices 7 to 14, so node 13 is at depth 3.
\end{solution}

\begin{solution}{exo:b1:heaps:4}
The first is not: node 3 holds 9 and its child at index 8 holds 7. The second is: 1 is at most 4 and 2, 4 at most 6 and
5, 2 at most 3.
\end{solution}

\begin{solution}{exo:b1:heaps:5}
The first extraction moves 4 to the root of $[4, 5, 2, 6, 7, 3]$; the left child 5 is not smaller, so the loop stops,
although the right child 2 is. The \omterm{def:b1:heaps:binary-heap}{heap property} is broken at the root, and the next extraction returns 4 instead of 2;
the errors then pile up, and 2 comes out sixth. Sift-down must choose the smaller of the two children. The smallest
heap that goes wrong in one extraction is $[1, 3, 2, 4]$: 4 moves to the root, is exchanged with 3 only, and gives $[3,
4, 2]$, whose root is not the minimum.
\end{solution}

\begin{solution}{exo:b1:heaps:6}
With 0-based indices, $i/2$ is the parent of $i$ only by chance. Node 0 is its own parent and its own left child. Node 2
is told that its parent is node 1 (key 3), which looks like a broken heap, but its parent is node 0. Node 4's parent is
node 1, not node 2. The left children $2i$ point one cell too early (node 1 to node 2, its own sibling). The correct
formulas are $\lfloor (i - 1)/2 \rfloor$ and $2i + 1$, or, keeping the 1-based ones, an \omterm{def:b1:pointers-and-arrays:array}{array} whose cell 0 is unused.
\end{solution}

\begin{solution}{exo:b1:heaps:7}
\omcodes{code/bachelor-1/18-heaps/exercises/gheap.c}{23-42}{Insertion and extraction of elements of any type}
Elements are moved with \code{memcpy}, \code{size} bytes at a time; the element at index $i$ starts at byte $i \cdot
\code{size}$. The comparison function returns a negative, zero or positive number, as for \code{qsort}; reversing it
turns the max-heap into a min-heap, which the test uses to serve the most urgent patient first.
\end{solution}

\begin{solution}{exo:b1:heaps:8}
\omcode{code/bachelor-1/18-heaps/exercises/kth.c}{4}{16}{The $k$-th smallest with a heap of size $k$}
The heap keeps the $k$ smallest values seen so far, negated, so that its root is the largest of them. A new value
smaller than that root replaces it. At the end the root is the $k$-th smallest. Each value costs $O(\log k)$; the test
compares with a sorted copy for many $n$ and $k$.
\end{solution}

\begin{solution}{exo:b1:heaps:9}
A \omterm{def:b1:heaps:complete}{complete} tree of height $h$ has its depths $0$ to $h - 1$ full, $2^h - 1$ nodes, and between 1 and $2^h$ nodes at depth
$h$: $2^h \leq n \leq 2^{h+1} - 1$. So $h \leq \log_2 n < h + 1$, and $h = \lfloor \log_2 n \rfloor$.
\end{solution}

\begin{solution}{exo:b1:heaps:10}
Let $S_k$ be the sum for the perfect tree with $2^k - 1$ nodes. $S_1 = 0 = 1 - 1$. The tree for $k + 1$ is a root of
height $k$ above two trees for $k$: $S_{k+1} = 2S_k + k = 2(2^k - 1 - k) + k = 2^{k+1} - 1 - (k + 1)$.
\end{solution}

\begin{solution}{exo:b1:heaps:11}
\omcode{code/bachelor-1/18-heaps/exercises/median.c}{11}{18}{Adding a number to the two halves}
Invariant: every value of the lower heap is at most every value of the upper heap, and the lower heap has as many
elements as the upper one or one more. Its root is then the lower median. A new value goes to the side its value
belongs to; if the sizes then differ by two, one root moves to the other heap, which keeps both properties. Each
addition makes at most three heap operations, $O(\log n)$. The test compares the median after every addition with a
sorted copy.
\end{solution}

\begin{solution}{exo:b1:heaps:12}
A sift-down from node $i$ makes at most $h_i$ iterations, each with at most two comparisons. \omterm{def:b1:heaps:sift}{Heapify} makes at most
$2\sum_i h_i = 2(n - \nu(n)) \leq 2n - 2$ comparisons (problem, Part IV). The extraction for a heap of $e$ elements, $e
= 1, \dots, n - 1$, sifts from the root of a tree of height $\lfloor \log_2 e \rfloor$: at most $2 \sum_{e=1}^{n-1}
\lfloor \log_2 e \rfloor$. With $L = \lfloor \log_2 (n - 1) \rfloor$, the value $j < L$ is taken $2^j$ times and $L$ is
taken $n - 2^L$ times, so the sum is $\sum_{j < L} j 2^j + L(n - 2^L) = (L - 2)2^L + 2 + Ln - L2^L = Ln - 2^{L+1} + 2 \leq
Ln - n + 2$, since $2^{L+1} \geq n$. In all, at most $2Ln - 2n + 4 + 2n - 2 = 2Ln + 2$ comparisons. Finally $\log_2 n - L
\geq \log_2 \frac{n}{n - 1} \geq \frac{1}{n \ln 2}$, so $2n \log_2 n \geq 2Ln + \frac{2}{\ln 2} > 2Ln + 2$. The
measurements give 1.70 to 1.85 times $n \log_2 n$.
\end{solution}

\begin{solution}{pb:b1:heaps:1}
\textbf{1.} Each output compares the heads of the non-empty lists, at most $k$ of them: $k - 1$ comparisons, and $(k -
1)n$ in all.

\textbf{2.} $31 \times 378 = 11\,718$ against 10\,884: towards the end lists run out, and fewer heads are compared.

\textbf{3.} With two lists both methods make one comparison per output while both lists last, and none after.

\textbf{4.} $\log_2 k$ rounds; in each round every element is in one merge, which makes at most one comparison per
element output: at most $n \log_2 k$ comparisons, like the heap.

\textbf{5.} The next output is always one of the heads, so the merge reads each list sequentially and holds $k$ values:
lists in files are read as streams, with memory for $k$ values only.

\textbf{6.} Each list is sorted, so every value not yet output is at least the head of its list, hence at least the
smallest head, which is the value output. So the outputs never decrease.

\textbf{7.} The heap holds exactly one entry for each non-empty list: the current head of that list.

\textbf{8.} The root is the smallest entry, the smallest head, which by question 6 is the smallest value not yet output.

\textbf{9.} The next head of the same list takes the root's place and is sifted down once. Extract-min followed by an
insertion would \omterm{def:b1:heaps:sift}{sift down} the last \omterm{def:b1:induction-and-termination:rooted-tree}{leaf}, then \omterm{def:b1:heaps:sift}{sift up} the new value: up to twice the work.

\textbf{10.} The heap has at most $k$ entries and height at most $\lfloor \log_2 k \rfloor$; a sift-down makes at most
two comparisons per level. In all at most $2n \lfloor \log_2 k \rfloor$, plus the $k$ initial insertions.

\textbf{11.} $2 \times 378 \times 5 = 3780$, and the table shows 2720, including the insertions.

\textbf{12.} The integer orders first by day, then by branch: the entries are distinct and the order of equal days is
fixed, and the branch is recovered by \texttt{\% 64}. It assumes $k \leq 64$ (and days small enough not to overflow an
\code{int}).

\textbf{13.} No: $k$ insertions cost $O(k \log k)$, against $O(n \log k)$ for the merge, with $k$ much smaller than $n$.

\textbf{14.} 8.01, 10.00, 12.00, \dots, 18.00 for $n = 2^{10}, 2^{12}, \dots, 2^{20}$ (minus 1): about $\log_2 n - 2$ per
key, since half the keys are leaves that climb the whole height, and the average depth is $\log_2 n - 2$ or so.

\textbf{15.} A node's key is larger than all the keys below it, so at each step it is larger than its smaller child and
goes down, until it has no child.

\textbf{16.} $1023 = 1111111111_2$, $\nu = 10$: $1013$. $1\,048\,575 = 2^{20} - 1$, $\nu = 20$: $1\,048\,555$. Both
match.

\textbf{17.} Each iteration compares the two children, then the node with the smaller one; an exchange follows, or the
sift stops. So a sift of $s$ exchanges makes at most $2(s + 1)$ comparisons, and \omterm{def:b1:heaps:sift}{heapify} at most $2(n - \nu(n)) + 2
\lfloor n/2 \rfloor < 3n$.

\textbf{18.} The $n - 1$ extractions, each sifting from the root: $\Theta(n \log n)$ in the worst case, against $O(n)$ for
\omterm{def:b1:heaps:sift}{heapify}.

\textbf{19.} The deepest descendants of a node include those on its leftmost path: going $d$ levels down to the left from
node $i$ reaches node $(i + 1)2^d - 1$. It exists if and only if $(i + 1)2^d \leq n$, so the height is the largest such
$d$, $\lfloor \log_2 (n/(i + 1)) \rfloor$.

\textbf{20.} $h_i \geq h$ exactly when $i + 1 \leq n/2^h$, for $\lfloor n/2^h \rfloor$ nodes; a node of height $h_i$ is
counted for $h = 1, \dots, h_i$, so $S(n) = \sum_i h_i = \sum_{h \geq 1} \lfloor n/2^h \rfloor$.

\textbf{21.} $\sum_{h=1}^{j} 2^{j-h} = 2^{j-1} + \dots + 1 = 2^j - 1$. Write $n = \sum_j 2^j$ over its bits: $n/2^h$ is
$\sum_{j \geq h} 2^{j-h}$, an integer, plus $\sum_{j < h} 2^j/2^h < 1$; so $\lfloor n/2^h \rfloor = \sum_{j \geq h}
2^{j-h} = \sum_j \lfloor 2^j/2^h \rfloor$.

\textbf{22.} Summing over $h$: $S(n) = \sum_{\text{bits } 2^j \text{ of } n} (2^j - 1) = n - \nu(n)$. For $n = 1000 =
1111101000_2$, $\nu = 6$ and $S = 994$, as printed.

\textbf{23.} $\nu(n) \leq \lfloor \log_2 n \rfloor + 1$: the bound $n$ is off by at most that much.

\textbf{24.} \emph{The heap identity}: in a heap of $n$ nodes the heights add up to $n - \nu(n)$ (questions 19 to 22),
checked by program for every $n \leq 10^4$. Hence \omterm{def:b1:heaps:sift}{heapify} makes fewer than $n$ exchanges and fewer than $3n$ comparisons,
and exactly $n - \nu(n)$ exchanges on decreasing keys, as \cref{out:b1:heaps:build} shows.
\end{solution}
